\documentclass[a4paper,11pt]{article}

\usepackage[T1]{fontenc}
\usepackage{lmodern}
\usepackage[a4paper,margin=2.7cm]{geometry}
\usepackage{amsmath,amssymb}
\setlength{\marginparwidth}{2cm}
\usepackage[theorems,framedtheorems]{tssuite}

\title{\texttt{tssuite} Theorem Examples}
\author{Focused example}
\date{\today}

\NewNumberedTheorem{theorem}{Theorem}
\NewNumberedTheorem{lemma}{Lemma}[sibling=theorem]
\NewStarredTheorem{remark}{Remark}
\NewDocumentTheorem{definition}{Definition}
\NewNumberedTheorem{exampletheorem}{Example}[style=definition,parent=section]

\MakeFramedTheorem{ftheorem}{theorem}
\MakeFramedTheorem{fdefinition}{definition}

\begin{document}

\maketitle

This document loads \texttt{tssuite} with theorem support and framed theorem
wrappers. The theorem helpers are thin wrappers around \texttt{keytheorems}.

\section{Basic Theorems}

\begin{theorem}[label=thm:transitive,note=Transitivity]
If $a=b$ and $b=c$, then $a=c$.
\end{theorem}

\begin{lemma}[note=Shared counter]
If $x>0$ and $y>0$, then $x+y>0$.
\end{lemma}

The theorem can be referenced with the ordinary label mechanism:
Theorem~\ref{thm:transitive} is labelled through a \texttt{keytheorems} option.

\begin{remark}
This unnumbered remark was created with \verb|\NewStarredTheorem|.
\end{remark}

\section{Paired Theorem Environments}

\verb|\NewDocumentTheorem| creates numbered and unnumbered variants from one
declaration.

\begin{definition}[note=Numbered]
A monoid is a set with an associative binary operation and an identity element.
\end{definition}

\begin{definition*}
This is the unnumbered variant created beside the numbered definition
environment.
\end{definition*}

\section{Framed Theorems}

Framed theorem wrappers place an existing theorem-like environment inside an
\texttt{mdframed} box.

\begin{ftheorem}[Cancellation]
If $a+c=b+c$ in an abelian group, then $a=b$.
\end{ftheorem}

\begin{fdefinition}[Group]
A group is a monoid in which every element has an inverse.
\end{fdefinition}

\section{Stored And Listed Theorems}

\begin{theorem}[store=stored-example,note=Stored theorem]
Every finite subset of a totally ordered set has a maximal element.
\end{theorem}

The stored theorem can be printed again later:

\getkeytheorem{stored-example}

\begin{exampletheorem}[note=Section-local numbering]
This theorem-like environment is numbered within the current section.
\end{exampletheorem}

\listofkeytheorems[title={Theorems in this focused example}]

\end{document}
